\section{Các mối đe dọa đến tính hợp
lệ}\label{sec:threats}

\subsection{Tính hợp lệ nội
tại}\label{sec:sec:validitynoitai}

Mối đe dọa chính đến tính hợp lệ nội tại là trạng thái compute không
được kiểm soát hoàn toàn. Thực nghiệm chạy trong môi trường Databricks
do cloud quản lý, trong đó cluster khởi động, trạng thái cache, cơ chế
lập lịch và phân bổ tài nguyên không cố định. Các yếu tố này có thể ảnh
hưởng đến overhead thời gian chạy độc lập với treatment. Để giảm rủi ro
này, các lần chạy khởi động (warm-up) được loại khỏi phân tích chính và
median được ưu tiên hơn mean.

Một mối đe dọa khác là hiệu ứng lịch sử. Nếu các bảng không được reset
đúng giữa các kịch bản tamper, dữ liệu cũ hoặc bản ghi ledger cũ có thể
tạo dương tính giả hoặc âm tính giả. Vì vậy, quy trình bao gồm bước
\texttt{RESET\_BASELINE} bắt buộc trước khi chuyển sang kịch bản tiếp theo.

\subsection{Tính hợp lệ ngoại
tại}\label{sec:validitingoaitai}

Thực nghiệm sử dụng một workspace Databricks và dữ liệu giao dịch tổng
hợp. Kết quả có thể không khái quát trực tiếp sang AWS Glue, Spark tại
chỗ, dbt, Snowflake hoặc môi trường lakehouse production. Dữ liệu tổng
hợp cũng thiếu các đặc điểm như phân phối lệch, trôi cấu trúc dữ liệu,
giá trị thiếu và độ phức tạp vận hành của dữ liệu thực tế. Các nghiên
cứu lặp lại trong tương lai nên triển khai trên nhiều môi trường và sử
dụng dữ liệu giống môi trường vận hành thực tế đã ẩn danh.

\subsection{Tính hợp lệ khái
niệm}\label{sec:validitykhainiem}

Hệ thống không nên được diễn giải là Blockchain đầy đủ. Khái niệm được
đánh giá trong nghiên cứu là Hash-linked Tamper-evident Ledger. Cơ chế
này cung cấp bằng chứng toàn vẹn thông qua các giá trị hash và liên kết
\texttt{previous\_hash}, nhưng không cung cấp đồng thuận phân tán hay
tính bất biến của public chain. Một mối đe dọa khái niệm khác là chỉ số
tỷ lệ phát hiện. Tỷ lệ phát hiện 100\% chỉ áp dụng cho sáu kịch bản
tamper được lựa chọn trong thực nghiệm.

\subsection{Tính hợp lệ kết
luận}\label{sec:validityketluan}

Số lượng kịch bản tamper, các mức số lượng bản ghi và số lần lặp lại đo
hiệu năng còn hạn chế. Điều này làm cho kết quả hữu ích như bằng chứng
bán thực nghiệm, nhưng chưa đủ cho kết luận thống kê mạnh. Các giá trị
overhead mang tính xấp xỉ vì được đọc từ dashboard quan sát thay vì xuất
trực tiếp từ dữ liệu benchmark thô. Cần ít nhất khoảng 30 lần chạy lặp
lại cho mỗi mức số lượng bản ghi để tính khoảng tin cậy, biểu đồ phân
phối và phân tích cỡ hiệu ứng mạnh hơn.
